\subsection*{Chapter 25: Instruction Fine-Tuning and Parameter-Efficient Adaptation}

\paragraph{Exercise 25.1.} Replace \code{format\_alpaca} by \code{format\_phi3} in both the
dataset and the generation prompt, and extract responses by splitting on
\code{<|assistant|>} instead of the Alpaca header. The Phi-3 template removes roughly
25--30 tokens per example, so total training tokens drop by roughly a quarter to a
third and epoch time falls accordingly. Judge scores are usually similar. Note that
\code{<|user|>} is not a special token for GPT-2 and is split into several BPE pieces,
which is harmless but slightly less compact than the name suggests.

\paragraph{Exercise 25.2.} Losses averaged over response tokens only are not comparable
with losses averaged over all tokens, because prompt tokens (especially the fixed
preamble) are easy to predict and pull the full-sequence average down. Evaluate both
models on the same response-only target set. Response-only training usually matches or
slightly beats full-sequence training on judge score. Full-sequence loss can win when
responses are very short and prompts carry useful in-domain text, because it then
provides much more gradient signal per example.

\paragraph{Exercise 25.3.} Two microbatches with $n_1=10$ and $n_2=90$ target tokens and
mean losses $\bar\ell_1,\bar\ell_2$: averaging the means gives
$\tfrac12(\bar\ell_1+\bar\ell_2)$, so each token of microbatch 1 has weight $1/20$ and
each token of microbatch 2 has weight $1/180$. The per-token mean of the full step gives
every token weight $1/100$. Fix: compute \code{F.cross\_entropy(..., reduction="sum")}
per microbatch, count non-ignored targets, and divide the accumulated sum by the total
count before \code{backward} (or scale each microbatch's sum by $1/n_{\text{total}}$,
which requires knowing $n_{\text{total}}$ in advance, for example by pre-counting the
step's targets).

\paragraph{Exercise 25.4.} Random batches with dynamic padding typically waste 30--50\% on
instruction data with a long-tailed length distribution; bucketing reduces this to a few
percent; FFD packing into 512- or 1{,}024-token rows leaves only the tail of the last
rows empty, usually below 2\%. For attention, a dense row mask costs $T^2$ per row
regardless of content. A varlen kernel costs $\sum_i\ell_i^2$, which for documents of
about 100 tokens packed into $T=1024$ is about 10 documents $\times\,100^2 = 10^5$
versus $1024^2\approx10^6$, a $10\times$ reduction in attention FLOPs.

\paragraph{Exercise 25.5.} Probe: fix document 2, vary document 1 in the same packed row,
and compare the logits at a position in document 2. With naive packing they change.
With the block-diagonal mask and position reset they are identical up to floating-point
error. After fine-tuning on naively packed data the model learns to use the preceding
documents, for example by copying formatting or answers from the previous example.
The effect is visible as degraded single-example performance at inference time, where
no preceding documents exist.

\paragraph{Exercise 25.6.} Reference answers should score near the top of the scale and
wrong answers near the bottom; deviations reveal judge bias or prompt problems.
Correlations between two different judges are typically high (Pearson 0.7--0.9) on
clearly separable answers but lower among mediocre ones, and mean absolute differences
of 5--15 points are common. A single reported score therefore has an uncertainty of
several points from the judge alone, on top of sampling uncertainty from 110 test
items (standard error $\approx\sigma/\sqrt{110}$).

\paragraph{Exercise 25.7.} With $\alpha=2r$, the effective update scale $\alpha/r=2$ is
constant, so the magnitude of $\Delta W$ produced by a given learning rate is roughly
independent of $r$, and one learning rate works across ranks. With fixed $\alpha=16$,
$\alpha/r$ shrinks as $r$ grows (16 at $r=1$, 0.25 at $r=64$), so large ranks
effectively train with a smaller step and appear worse unless the learning rate is
re-tuned. Quality usually saturates by $r=8$--$16$ for a task of this size.

\paragraph{Exercise 25.8.} For a fixed parameter budget, adapting all linear layers
(with a smaller $r$) usually wins on SFT, because MLP layers hold much of the
task-relevant capacity. The spam task is easy enough that most placements reach
similar accuracy. Attention-only placements are slightly weaker on SFT.

\paragraph{Exercise 25.9.} In full fine-tuning the weight-gradient GEMMs are about a third
of matmul time. With LoRA they disappear for frozen weights, while the adapter GEMMs
are tiny but numerous (two extra small matmuls per adapted layer in each direction).
On small models these extra kernels and the optimizer step can eat most of the
theoretical $1.5\times$ saving, so measured speed-ups are often smaller than predicted.
On large models the prediction is approached.

\paragraph{Exercise 25.10.} Block-normalised GPT-2 weights are roughly bell-shaped but
heavier-tailed than a Gaussian, with occasional outlier channels. NF4 gives lower
reconstruction error than uniform int4 for the bulk of the weights, because its
levels are denser near zero, and typically a smaller perplexity increase. Uniform int4
does relatively better on blocks dominated by an outlier, where the absmax scale
squeezes the remaining weights into a few levels for both schemes.

\paragraph{Exercise 25.11.} With one adapter, the gather is a broadcast and the cost is
close to the base model plus two small matmuls. With 8 or 64 distinct adapters, the
gather materialises per-row copies of $\mat{A}$ and $\mat{B}$ ($B\times r\times d$
elements), which adds memory traffic proportional to the batch size. That is where
throughput drops. Fused kernels (Punica's BGMV/SGMV) read each row's adapter directly
from the adapter pool inside the matmul, never materialising the gathered copies, and
group rows that share an adapter (segmented GEMM) to recover tensor-core efficiency.
